
\addcontentsline{toc}{chapter}{符号与约定}
\begin{center}
  {\bfseries\Large 符号与约定}
\end{center}

\bigskip\noindent
本书在全文中统一采用如下记号约定；各章给出定义时不另立与此冲突的记号。

% ============================================================
\section*{一、主要函数与符号一览}

\renewcommand{\arraystretch}{1.7}
\begin{longtable}{@{}lp{10cm}@{}}
\toprule
\textbf{符号} & \textbf{含义} \\
\midrule
$\# S$ 或 $\#\{\dots\}$
  & 集合 $S$ 或满足括号内条件的集合的元素个数 \\
$\pi(x)$
  & 不超过 $x$ 的素数个数：
    $\pi(x) = \#\{p \le x : p \text{ 为素数}\}
    =\displaystyle\sum_{p\le x,\,p\text{ 素数}}1$
    （右连续阶梯函数） \\
$\pi_0(x)$
  & 素数计数函数在跳跃点处的正则化版本：
    $\pi_0(x)=\dfrac{\pi(x^-)+\pi(x^+)}{2}$，
    即在素数 $p$ 处取左右极限的平均值，$\pi_0(p)=\pi(p)-\dfrac{1}{2}$；
    在非素数处 $\pi_0(x)=\pi(x)$。
    显式公式在素数处给出的是 $\pi_0$ 的值，而非 $\pi$ 的值 \\
$J(x)$
  & 黎曼素数幂计数函数，$J(x)=\displaystyle\sum_{p^k\le x}\dfrac{1}{k}$（右连续） \\
$J_0(x)$
  & $J$ 在跳跃点的左右平均：$J_0(x)=\dfrac{J(x^-)+J(x^+)}{2}$；
    非素数幂处 $J_0(x)=J(x)$ \\
$\theta(x)$
  & 第一 Chebyshev 函数，$\theta(x)=\displaystyle\sum_{p\le x}\log p$ \\
$\psi(x)$
  & 第二 Chebyshev 函数，$\psi(x)=\displaystyle\sum_{n\le x}\Lambda(n)$（右连续） \\
$\psi_0(x)$
  & $\psi$ 在跳跃点的左右平均：$\psi_0(x)=\dfrac{\psi(x^-)+\psi(x^+)}{2}$；
    非素数幂处 $\psi_0(x)=\psi(x)$。
    \textbf{显式公式}（Riemann--von Mangoldt 型）在素数幂处给出的是 $\psi_0$ 的值 \\
$\mu(n)$
  & Möbius 函数：$n$ 含平方因子时取 $0$；无平方因子时取 $(-1)^k$（$k$ 为不同素因子个数） \\
$M(x)$
  & Mertens 函数，$M(x)=\displaystyle\sum_{n\le x}\mu(n)$ \\
$\Lambda(n)$
  & von Mangoldt 函数：$\Lambda(n)=\log p$（当 $n=p^k$），否则 $\Lambda(n)=0$ \\
$\Li(x)$
  & \textbf{偏移对数积分}：实函数
    $\Li:(1,+\infty)\to\mathbb{R}$，
    $\Li(x)=\displaystyle\int_2^x\dfrac{\mathrm{d}t}{\log t}$。
    素数定理、von Koch、$\pi\sim\Li$ 用此。
    定义域不含复自变量；复点用 $\li(z)$。
    见第~\ref{ch:logint_Ei}~章第~\ref{sec:Li}~节 \\
$\li(x)$，$\li(x^{\rho})$
  & \textbf{对数积分}：实轴 Cauchy 主值；复幂路径主值/延拓。
    显式公式用 $\li(x)$ 与 $\li(x^{\rho})$。
    实轴 $\Li=\li-\li(2)$ 仅用于实比较互译。
    见第~\ref{ch:logint_Ei}~章第~\ref{sec:Li-li-convention}~、\ref{sec:li-complex-power}~节 \\
$\Ei(z)$
  & 指数积分。$\li(x)=\Ei(\log x)$，$\li(x^{\rho})\leftrightarrow\Ei(\rho\log x)$
    （第~\ref{ch:logint_Ei}~章第~\ref{sec:Ei}~节；算法见数值章） \\
$\zeta(s)$
  & 黎曼 $\zeta$ 函数；在 $\Real(s)>1$ 时由 Dirichlet 级数与 Euler 乘积定义，并解析延拓至 $\mathbb{C}\setminus\{1\}$ \\
$\Gamma(s)$
  & Gamma 函数；在 $\Real(s)>0$ 时由 Euler 积分
    $\Gamma(s)=\displaystyle\int_0^\infty t^{s-1}e^{-t}\,\mathrm{d}t$ 给出，
    并解析延拓为 $\mathbb{C}$ 上的亚纯函数（极点在非正整数）。
    系统讨论见第~\ref{ch:gamma_function}~章 \\
$\xi(s)$
  & 完全 $\xi$ 函数（整函数），
    $\xi(s)=\dfrac{1}{2}s(s-1)\pi^{-s/2}\Gamma\!\left(\dfrac{s}{2}\right)\zeta(s)$ \\
$\rho=\beta+i\gamma$
  & $\zeta(s)$ 的非平凡零点，$0<\beta<1$，$\gamma\in\mathbb{R}$ \\
$N(T)$
  & 虚部在 $(0,T]$ 内的非平凡零点个数（计重数） \\
$N_0(T)$
  & 其中位于临界线 $\Real(s)=\dfrac{1}{2}$ 上的个数；RH $\iff N_0(T)=N(T)\ (\forall T)$ \\
$\vartheta(t)$
  & \textbf{Riemann--Siegel 相位}（Hardy 理论）：
    $\vartheta(t)=\Imag\log\Gamma\bigl(\dfrac{1}{4}+\dfrac{it}{2}\bigr)-\dfrac{t}{2}\log\pi$
    （与 Chebyshev $\theta(x)$、Jacobi $\vartheta$ 不同；分工见第三节） \\
$Z(t)$
  & \textbf{Hardy $Z$ 函数}：$Z(t)=e^{i\vartheta(t)}\zeta\bigl(\dfrac{1}{2}+it\bigr)\in\mathbb{R}$；
    $\bigl|Z(t)\bigr|=\bigl|\zeta\bigl(\dfrac{1}{2}+it\bigr)\bigr|$；
    极坐标形式 $\zeta\bigl(\dfrac{1}{2}+it\bigr)=Z(t)e^{-i\vartheta(t)}$。
    用于临界线零点数值计算与几何示意，\textbf{不是}新的 $L$-函数 \\
$\mathrm{RH}$
  & 黎曼猜想（RH，Riemann Hypothesis）：$\zeta(s)$ 的所有非平凡零点满足 $\beta=\dfrac{1}{2}$ \\
$\Real(s),\;\Imag(s)$
  & 复数 $s$ 的实部与虚部（亦写作 $\Re(s),\Im(s)$） \\
$\mathbb{N},\mathbb{Z},\mathbb{Q},\mathbb{R},\mathbb{C}$
  & 自然数、整数、有理数、实数、复数集 \\
\bottomrule
\end{longtable}

\paragraph{主要函数的书写（$\zeta$、$\Gamma$、$\xi$ 等）}
为避免「$\zeta$ / $\zeta(s)$ / $\zeta(s)$ 函数」混用，全文统一如下（$\Gamma$、$\xi$ 同理；
$\pi(x)$、$J(x)$、$\psi(x)$ 等算术计数函数始终带自变量写法，不在此列）。

\begin{enumerate}[leftmargin=2em,itemsep=0.35em]
  \item \textbf{表达式与定义}（出现自变量或等号时）：写 $\zeta(s)$、$\Gamma(s)$、$\xi(s)$。\\
        例：$\zeta(s)=\sum_{n=1}^\infty n^{-s}$；$\xi(s)=\xi(1-s)$；
        非平凡零点满足 $\operatorname{Re}s=\frac{1}{2}$。
  \item \textbf{叙述中指整个函数}（不强调某一取值）：\\
        该章（或全书）\textbf{第一次}出现时，写「$\zeta$ 函数」「$\Gamma$ 函数」「$\xi$ 函数」
        （$\zeta$ 亦可写「黎曼 $\zeta$ 函数」）；\\
        \textbf{此后}同一叙述中可只写 $\zeta$、$\Gamma$、$\xi$。\\
        例：首次「$\zeta$ 函数的 Euler 乘积」；后文「$\zeta$ 的非平凡零点」「$N(T)$ 计数 $\zeta$ 的零点」。
  \item \textbf{禁止冗余}：不写「$\zeta(s)$ 函数」「$\Gamma(s)$ 函数」「$\xi(s)$ 函数」
        （已有自变量，不必再加「函数」二字）。
  \item \textbf{定义句可稍长}（仅在给出定义或写出积分/级数表示时）：\\
        允许「黎曼 $\zeta$ 函数 $\zeta(s)$」「$\Gamma$ 函数（Euler 积分）
        $\Gamma(s)=\int_0^\infty t^{s-1}e^{-t}\,\mathrm{d}t$」这类一次写全的句式；\\
        定义之后仍按第 1--2 条：式子用 $\zeta(s)$，叙述用「$\zeta$ 函数」或 $\zeta$。
\end{enumerate}

\paragraph{西文人名与定理专名（一律拉丁拼写）}
作为数学专著，\textbf{西数学家人名与以其命名的定理/公式专名一律用拉丁拼写}，
不采用「哈达玛」「勒让德」「柯西」等中译作专名，也不写「哈达玛（Hadamard）」式夹注。
与英文文献、索引及国际通行写法一致。

\begin{enumerate}[leftmargin=2em,itemsep=0.35em]
  \item \textbf{人名}：Riemann, Euler, Gauss, Legendre, Dirichlet, Hadamard,
        Weierstrass, Mertens, Cauchy, Jensen, von Mangoldt, Perron, Hardy, Siegel 等。\\
        叙述例：Riemann（1859）；Hadamard 与 de~la~Vallée-Poussin（1896）。
  \item \textbf{专名（人名作定语）}：Euler 乘积、Dirichlet 级数、Hadamard 乘积、
        Riemann--Siegel 公式、Mertens 定理、von Mangoldt 函数、Perron 公式、
        Cauchy--Riemann 方程、Legendre--Gauss 型主项等。
  \item \textbf{连接号}：两名人名并列用一字线（en-dash）：
        Riemann--Siegel、Cauchy--Riemann、Legendre--Gauss、Wilbraham--Gibbs。
  \item \textbf{中文「黎曼」仅限固化术语}：\\
        叙述主语与年代写 Riemann（1859），不写「黎曼于 1859 年……」。\\
        中文「黎曼」\textbf{仅}用于下列已固化的术语：
        \begin{itemize}
          \item 黎曼猜想（书名、章题及 RH 的中文标准称呼）；
          \item 黎曼 $\zeta$ 函数、黎曼显式公式等。
        \end{itemize}
\end{enumerate}

\paragraph{版式：定理环境与框}
结构靠\textbf{编号与标题文字}区分；描边色统一为三类（结构 / 辅助 / 功能；方案 II，
见 \texttt{preamble.tex}）。\textbf{证明无框}；章内不另造新色框。

\begin{enumerate}[leftmargin=2em,itemsep=0.35em]
  \item \textbf{结构框}（冷蓝灰细线描边，同色）：定义、定理、引理、命题、推论、算法；
        以标题「定义 / 定理 / …」与节内编号区分类型。
  \item \textbf{辅助框}（中性灰细线，线略轻）：例、注。
  \item \textbf{功能框}（深灰细线）：主要用于原稿「解读」
        （第~\ref{ch:riemann_original_interpretation}~章）；
        另有要点框、结果框、说明框（同为细线、无填充），章内按需使用。
  \item \textbf{框的外观}：四周细线描边，框内\textbf{无填充底色}；
        结构/辅助框标题写在框内起首，加粗黑字（如「定义~2.1」「定理~3.2」）。
  \item \textbf{公式}：以编号方程为主；少数核心恒等式可加框强调。
  \item \textbf{公式体量（\texttt{\textbackslash displaystyle} / 分式）}：
        \begin{itemize}
          \item \textbf{独立公式}（\texttt{equation} / \texttt{align} /
                \texttt{\textbackslash[...\textbackslash]}）：
                环境本身已是显示型，\textbf{不必}再写 \texttt{\textbackslash displaystyle}；
                分式用 \texttt{\textbackslash frac} 即可（显示型下与 \texttt{\textbackslash dfrac} 同大）。
          \item \textbf{正文行内}：\textbf{不加} \texttt{\textbackslash displaystyle}；
                分式用 \texttt{\textbackslash frac} 或简写 $a/b$（如 $1/z$、$f'/f$）。
                上下标多、易挤时，优先改为独立公式，勿在行内强行 display。
          \item \textbf{表、章末汇总框（\texttt{\textbackslash fbox} 等）、本节符号一览}：
                为醒目可用 \texttt{\textbackslash displaystyle} 与 \texttt{\textbackslash dfrac}。
        \end{itemize}
  \item \textbf{章末}：各章宜有「本章小结与后续展望」，含本章要点与对后文的衔接；
        末章可写开放问题与延伸阅读。
        可选「重要定理与核心公式汇总」表（现多用 \texttt{\textbackslash fbox} 表，亦可用结果框），
        置于小结前或小结中。
\end{enumerate}

% ============================================================
\section*{二、跳跃点与 \texorpdfstring{$\psi_0,J_0,\pi_0$}{psi0,J0,pi0}}

未加下标的 $\psi,J,\pi$ 默认\textbf{右连续}。
显式公式与 Perron 公式在跳跃点给出左右平均，严格写法为 $\psi_0$、$J_0$、$\pi_0$
（定义见上表）。
声明「$x$ 非素数幂」或「非素数」时，$\psi_0=\psi$ 等，可省写下标。

% ============================================================
\section*{三、其他约定}

\begin{itemize}
  \item \textbf{对数}：本书中 $\log$ 均表示自然对数（以 $e$ 为底）。
        一般\textbf{不使用} $\ln$ 记号；凡涉及以 $10$ 为底的常用对数，一律写明 $\log_{10}$。
        这与解析数论国际文献（Titchmarsh、Edwards、Davenport 等）的惯例一致。
  \item \textbf{三种 “theta”}（全文统一，第~\ref{ch:riemann_numerical_approximation}~章详述）：
        $\theta(x)$ = 第一 Chebyshev 函数；
        $\vartheta(t)$ = Riemann--Siegel 相位（Hardy $Z$ 配套）；
        Jacobi $\vartheta$（模反演，见解析延拓章）与前二者均不同。
        Riemann--Siegel 相位只用 $\vartheta(t)$，不用 $\theta(t)$。
  \item $p$ 始终表示\textbf{素数}；$p^k$ 表示素数幂（$k\ge 1$）。
  \item $\sum_\rho$ 表示遍历 $\zeta(s)$ 所有非平凡零点的求和，
        按 $|\Imag(\rho)|$ 递增顺序\textbf{对称配对}（$\rho$ 与 $\bar\rho$ 同时加入），
        以保证条件收敛。
  \item 显式公式中的 Perron 围道积分取 $c>1$（$c$ 充分接近 $1$）。
  \item \textbf{渐近符号}（默认 $x\to\infty$；比较函数 $g$ 在所论范围内最终非零）：
        \begin{itemize}
          \item $f=O(g)$：存在常数 $C>0$，使对充分大的 $x$ 有 $|f(x)|\le C\,|g(x)|$
                （上界；$f$ 的增长不超过 $g$ 的常数倍）；
          \item $f=o(g)$：$f(x)/g(x)\to 0$
                （$f$ 比 $g$ 严格更小阶）；
          \item $f=\Omega(g)$：$\limsup_{x\to\infty}\bigl|f(x)/g(x)\bigr|>0$
                （$f$ 不是 $o(g)$；下界意义下「至少达到 $g$ 的阶」）。
                例：第~\ref{ch:elementary_number_theoretic_functions}~章由 Chebyshev 界得
                $\pi(x)=O(x/\log x)$，且存在 $c>0$ 使 $\pi(x)\ge c\,x/\log x$（更强的下界），
                因而亦有 $\pi(x)=\Omega(x/\log x)$；
          \item $f\sim g$：$f(x)/g(x)\to 1$（渐近等价，主项相同）。
        \end{itemize}
  \item 围道积分方向约定为\textbf{正向}（逆时针）。
\end{itemize}
